\mode*
\section{Method}
\label{method}

% TODO Anonymise the course (and institution) before submission; keep the
% concrete details for now.
We will study the students of prgi26, the autumn 2026 instance of the
introductory programming course DD1317 at KTH Royal Institute of Technology.
The course teaches Python to first-year students and is organised around a
series of laboratory assignments---from a first \enquote{Hello World} lab,
through functions and variables, input and error handling, repetitions and
lists, classes and objects, containers, and file handling---followed by an
individual project.
The course design of prgi26 will be based on the findings of the companion
paper \autocite{VTProgMisconceptions}: the teaching material is structured
around the patterns of variation derived there.
% XXX Update this paragraph when prgi26 is finalised in Canvas; the design is
% based on prgi25 with the variation-theory patterns from the companion paper.
This matters for the present study in two ways: the students will have
\emph{experienced} the patterns of variation as learners, and by the
generative-learning effect (\cref{debugging-as-learning}) we expect some of
them to start generating such patterns for themselves.

Not all students use the course material, however: some study on their own,
from other sources.
These students form a natural comparison group \emph{within} the cohort: they
do the same labs, but without experiencing the material's patterns of
variation.
We therefore also record which students use the course material, through the
learning platform's activity data complemented by self-report in the
questionnaire (\cref{material-use}).
We also have the students who know programming from before; we identify
them with a diagnostic quiz at the start of the course
(\cref{prior-knowledge}).

In summary, the method is as follows.
We record the students' every edit--run cycle with learnlog throughout the
course (\cref{learnlog-data}) and code the recorded debugging episodes for
patterns of variation (\cref{analysis}); we measure the approaches to
learning with the R-SPQ-2F (\cref{approach-measure}); and we relate
patterns, approach and outcomes to one another.
The questions are answered in two modes: descriptively
(\cref{rq:patterns,rq:experts}) and by attempting to falsify the
corresponding hypothesis (\cref{rq:deep,rq:success,rq:generative}).
Restating the questions and the hypotheses, the overview is as follows;
the details follow in the coming sections.

\rqpatterns*
Answered descriptively: coding the debugging episodes recorded with
learnlog (\cref{learnlog-data}) in variation-theoretic terms
(\cref{analysis}) yields the catalogue of which patterns appear, in what
order, and how often.
\Cref{h:patterns} additionally predicts:
\hpatterns*
As stated this is a universal claim, which a single counterexample would
falsify; since the coding is necessarily noisy, we operationalise it
statistically: pattern completeness should separate successful from
unsuccessful episodes.

\rqdeep*
Answered by testing \cref{h:deep}:
\hdeep*
The questionnaire (\cref{approach-measure}) provides the approach measure,
which we correlate with the coded patterns; the think-aloud subsample
(\cref{think-aloud}) validates both the coding and the classification.

\rqsuccess*
Answered by testing \cref{h:success}:
\hsuccess*
The debugging outcomes (\cref{outcomes}) provide the success measures; the
diagnostic quiz (\cref{prior-knowledge}) provides the control for prior
experience.

\rqexperts*
Answered descriptively---it has no hypothesis to test: we code the
teachers' and assistants' episodes identically (\cref{teacher-baseline})
and compare the two groups' patterns.

\rqgenerative*
Answered by testing \cref{h:generative}:
\hgenerative*
The material-use measure (\cref{material-use}) identifies the groups;
approach and prior experience are controlled for; since the groups are
self-selected, the answer is correlational.
Note, finally, that the overarching hypothesis---that debugging skill and
the deep approach to learning are the same underlying ability---is not
itself directly tested: \cref{h:deep,h:success} are consistent with the
two being identical, but also with two distinct abilities sharing a cause.
Answering the questions thus supports the overarching hypothesis by
inference to the best explanation, not by falsification.
The data collection follows in \cref{data-collection}, the analysis in
\cref{analysis}, the pilot of the coding in \cref{method-pilot}, and the
ethical considerations in \cref{ethics}.

\mode<presentation>{%
\begin{frame}
  \begin{block}{Setting}
    \begin{itemize}
      \item prgi26: introductory Python programming (DD1317, KTH), autumn
        2026.
      \item Labs 0--6 plus an individual project.
      \item Course design based on the companion paper's patterns of
        variation.
      \item Some students study on their own instead \(\to\) within-cohort
        comparison group.
    \end{itemize}
  \end{block}
\end{frame}%
}

\subsection{Data collection}
\label{data-collection}

We describe the data collection in the chronological order of the course:
what we collect at the start of the course
(\cref{prior-knowledge,approach-measure}), what we collect throughout it
(\cref{learnlog-data,think-aloud,material-use,teacher-baseline}), and what
derives from the recordings afterwards (\cref{outcomes}).

\subsubsection{Prior programming knowledge}
\label{prior-knowledge}

Some students know programming from before the course.
They matter twice: prior experience is a confound for \cref{h:success}, and
the experienced students form yet another group within the cohort, whose
debugging may exhibit the patterns from the start.
We therefore administer a diagnostic test---a quiz---at the start of the
course, measuring prior programming knowledge, complemented by
self-reported programming experience.
Besides prior programming knowledge, the start-of-course quiz captures the
following.
\begin{enumerate}
  \item Prior debugging habits, through an open question such as
    \enquote{when your program misbehaves, what do you do?}.
    The answers can themselves be coded for variation-thinking, giving a
    baseline debugging-approach measure that is independent of both the
    logs and the questionnaire---a third, independent angle on the
    connection between approach and patterns (\cref{rq:deep}).
  \item Languages known and years of experience, and whether they know
    Python specifically: experienced students will debug faster regardless
    of approach, so this is the control variable for \cref{h:success}.
  \item How they learned to program: in a formal course or
    self-taught---the self-taught being self-directed learners by
    definition, which bears directly on the deep approach
    (\cref{rq:deep}): they may already generate their own variation.
  \item Programming self-efficacy: less confident students may give up or
    persist differently, which confounds the persistence-based success
    measures of \cref{h:success}, the time and the number of cycles to
    fix.
  \item How they plan to study in this course, as a baseline for the
    material-use self-report (\cref{material-use}) and thereby for the
    generative-learning comparison it supports.
\end{enumerate}
For the knowledge part itself, a subset of a validated concept inventory
for introductory programming would make the measure comparable to other
studies.
% XXX Pick and verify a concept inventory (candidates to be reviewed)
% before finalising the quiz.

\subsubsection{Measuring the approach to learning: R-SPQ-2F}
\label{approach-measure}

We measure the students' approaches to learning with the revised two-factor
Study Process Questionnaire (R-SPQ-2F), which yields a deep-approach and a
surface-approach score from ten items each \autocite{RSPQ2F}.
Since the approach is not a fixed trait---students adapt it to their
conception of the task \autocite{MartonSaljo1976II}---we administer the
questionnaire twice, at the start and at the end of the course.
At the start of the course the students cannot yet answer with respect to
this course---they must have started the course for that---so the first
administration asks about their studies in general, to get a feel for their
general approach.
The second administration, at the end, asks with respect to this course;
those scores we treat as measurements of the approach in the context of
this course rather than as a stable property of the student.

\paragraph{Classifying the students compatibly with the original studies}

Some care is needed for our classification of students to be compatible with
the studies our theory rests on.
In the original studies the approaches were identified phenomenographically
and per task:
\textcquote[p.~142]{NCOL}{\textins{t}he distinction was based on the
learners' own accounts of how they went about learning (i.e., how they
experienced their own attempts to learn)}---the students were interviewed
about how they set about reading a particular text
\autocite{MartonSaljo1976I}.
The R-SPQ-2F is instead a self-report questionnaire; it does not classify an
actual learning episode but the student's reported way of studying.
We bridge this gap in three ways.
\begin{enumerate}
  \item In the end-of-course administration (\cref{approach-measure}) we
    have the students answer the R-SPQ-2F with respect to \emph{this
    course}, in line with the instrument's intended use by teachers
    evaluating the approaches of their own students \autocite{RSPQ2F} and
    with the task-dependence of the approach \autocite{MartonSaljo1976II}.
    The resulting scores are comparable with the questionnaire-based
    literature.
  \item In the think-aloud subsample (\cref{think-aloud}) we additionally
    classify the students the way \textcite{MartonSaljo1976I} did: we ask them to recount how they
    went about the lab and the debugging, and classify the accounts as deep
    or surface.
    This classification is the one compatible with \citeauthor{NCOL}'s
    variation-theoretic reinterpretation, and we report the agreement between
    it and the questionnaire classification.
  \item We guard against circularity.
    \Citeauthor{NCOL}'s reinterpretation \emph{defines} the deep approach as
    creating the necessary patterns of variation
    \autocite[p.~145]{NCOL}; if we were to classify the students' approaches
    from the same edit--run logs that we code for patterns of variation,
    \cref{h:deep} would hold by construction.
    The approach classification therefore uses only the questionnaire and the
    interview accounts, never the logs.
\end{enumerate}

\subsubsection{Recording the students' programming: learnlog}
\label{learnlog-data}

We record the students' programming with the \texttt{learnlog} tool
\autocite{learnlog}.
The students run \texttt{learnlog init} once in their course directory (one
directory for the whole course, not one per assignment or project); the
tool then installs itself into the directory's virtual environment and
activates on every Python run, so that every program run is recorded
transparently---even runs that fail before the program starts, such as
syntax errors: source code changes, command-line arguments, standard
input/output/error, and unhandled exceptions are all stored.
The data is stored in a hidden local Git repository, which gives a complete
timeline of the students' edit--run cycles.
The students push their logs to per-student Git repositories (or export them
as bundle files and submit them through the course platform), and tag the
events per lab so that episodes can be attributed to assignments.

\begin{frame}<presentation>[fragile]
  \begin{block}{learnlog\only<presentation>{~\autocite{learnlog}}}
    \begin{itemize}
      \item \texttt{learnlog init} once per course directory.
      \item Records every run: source diff, arguments, I/O, exceptions.
      \item Hidden local Git repo \(\to\) complete edit--run timeline.
      \item Students push logs or submit bundles; events tagged per lab.
    \end{itemize}
  \end{block}
\end{frame}

For the debugger, an edit--run cycle is exactly one experiment: a variation of
the program followed by an observation of its behaviour.
The sequence of edit--run cycles is therefore precisely the data in which
patterns of variation---what was varied, what was kept invariant, and in
what order---should be visible, if our hypotheses hold.
Recall the phenomenographic stance
(\cref{background}): acts express understanding, so the recorded acts are the
appropriate data for a second-order analysis.

\subsubsection{Think-aloud sessions with a subsample}
\label{think-aloud}

With a smaller subsample of student volunteers we also hold think-aloud
debugging sessions during the course, of the kind piloted in \cref{pilot}:
the student debugs while speaking their thoughts aloud, and we record the
session with learnlog (\cref{learnlog-data}) together with the audio and
video (probably over Zoom).
The subsample serves two purposes.
First, the logs record acts but not the reasoning between them, so the
think-aloud sessions let us validate that the patterns coded from the logs
(\cref{analysis}) correspond to the patterns in the students' reasoning.
Second, the sessions give us the students' task-specific accounts of how
they went about the debugging, which we need to classify their approaches
the way the original studies did (\cref{approach-measure}).

\subsubsection{Course material use}
\label{material-use}

We need to know who uses the course material to be able to answer
\cref{rq:generative}: the groups compared there are defined by whether the
student has experienced the material's patterns of variation.
To identify these groups we measure each student's use
of the course material: the learning platform (Canvas) records which pages,
modules and embedded activities each student has worked through, and we
complement this with a self-report item alongside the R-SPQ-2F asking how the
student mainly studies (from the course material, from other sources, or a
mix).
The interactive material assignments (FeedbackFruits) give us detailed data
about the students' interactions with the material as well.
% XXX Decide the exact operationalisation (threshold on activity data vs
% self-report) once the prgi26 Canvas course is finalised.

\subsubsection{The teachers' debugging: an expert baseline}
\label{teacher-baseline}

\Cref{rq:experts} asks about expert debuggers.
We cannot sample expert debuggers in general, so we use the course's teachers
and teaching assistants as our approximation of experts, and collect the same
kind of data from them as from the students.
They run learnlog while preparing course material and during live-coding
sessions (learnlog is designed for recording and replaying such sessions
\autocite{learnlog}), and we complement this with think-aloud debugging
sessions of the kind piloted in \cref{pilot}.
The approximation is imperfect---the teaching assistants in particular vary
in experience---so we record their programming experience alongside, which
also lets us relate pattern completeness to expertise within the group.
Their episodes are coded identically to the students'
(\cref{episode-coding}), giving an expert baseline: which patterns experts
generate, in what order, and how complete their episodes are.

\subsubsection{Debugging outcomes}
\label{outcomes}

For each debugging episode in the recordings (\cref{learnlog-data}) we
record whether the bug was fixed, the number of edit--run cycles, and the
wall-clock time from the first failing run to the first passing run.
These constitute our measures of debugging success for \cref{h:success}.

\subsection{Analysis}
\label{analysis}

We define a \emph{debugging episode} as the sequence of edit--run cycles from
a run that surprises the student (typically a failing run, an unhandled
exception, or wrong output) to the run where the behaviour is understood and
fixed.
The \texttt{learnlog analyse} command generates reports of the edit--run
cycles \autocite{learnlog}; on these we perform a qualitative,
variation-theoretic coding:\label{episode-coding}

\begin{frame}
\begin{enumerate}
  \item For each consecutive pair of runs in an episode: what did the student
    vary, and what did they keep invariant?
  \item Classify the steps into the patterns: \emph{contrast} (varying the
    suspected critical aspect against the surprising behaviour),
    \emph{generalisation} (varying non-critical aspects, keeping the suspected
    aspect invariant---the experiments), and \emph{fusion} (varying the
    aspects together, reconciling the new understanding with the old).
  \item Rate the episode's pattern completeness: which of the patterns are
    present, and whether the episode ends in a fix.
\end{enumerate}
\end{frame}

Note that when programming, the generalisation and fusion experiments are
probably not made in the lab file itself: they are typically done in a REPL
or a throw-away file.
Since learnlog activates on every Python run in the course directory's
virtual environment---including REPL sessions and scratch files---these
runs are recorded too, and the coding must consider all recorded runs, not
only those of the lab solution.

One possibility is that a student does not generate the generalisation and
fusion patterns at all, but instead searches for an explanation online or
asks an AI assistant (\cref{debugging-as-learning}).
In the logs such episodes show up as a sudden fix without preceding
experiments, possibly with pasted-in code; we code these as externally
sourced.
This matters for \cref{h:patterns}: a successful episode without patterns
is not a counterexample if the variation was outsourced---the hypothesis
concerns the episodes where the student generated the explanation
themselves.

This coding answers \cref{rq:patterns} (which patterns appear).
For \cref{rq:deep,rq:success} we relate the coded episodes to the R-SPQ-2F
scores and the debugging outcomes: correlating deep-approach scores with
pattern counts and completeness (\cref{h:deep}), and both with debugging
success (\cref{h:success}).

The within-cohort comparison group (\cref{material-use}) lets us answer
\cref{rq:generative} already in this study: if experiencing patterns of
variation breeds generating them, as \cref{h:generative} predicts, then the
students who use the course material should generate more, and more complete,
patterns in their debugging than the students who study on their
own---controlling for approach (R-SPQ-2F) and prior programming experience.
The groups are self-selected, not randomised, so this comparison is
correlational; we discuss the resulting confounds in \cref{discussion}.

For \cref{rq:experts} we compare the students' coded episodes with the
teachers' expert baseline (\cref{teacher-baseline}): which patterns each group generates, in
what order, and with what completeness.
If the overarching hypothesis is right, the experts' episodes should look
like the deep-approach students' episodes, only more so---and the
differences between the two should point out precisely what debugging
instruction needs to provide.

\subsection{Pilot}
\label{method-pilot}

As reported in \cref{pilot}, we have piloted the variation-theoretic coding on
think-aloud recordings of experienced programmers debugging the
default-argument phenomenon of \cref{debugging}.
The patterns were identifiable in the recordings, and the successful episodes
contained all of them.
The pilot used screen and voice recordings; the main study instead uses
learnlog's edit--run data, which scales to a full course but does not capture
the students' reasoning between runs.
We therefore complement the logs with the think-aloud subsample
(\cref{think-aloud}), to validate that the patterns coded from the logs
correspond to the patterns in the students' reasoning.

\subsection{Ethics}
\label{ethics}

Participation in the study is opt-in and requires informed consent; the
students' grades are unaffected by participation, and the logs are collected
and analysed pseudonymously.
The learnlog data contains the students' source code and program
interactions, but no more than what the students already submit with their
lab solutions.
% XXX Ethical review application for prgi26 data collection: reference number
% and details to be added.
